\properlane{prayer-and-conversion}{The confident petitioner asks to be changed}

The Communion promises an answer to whatever is asked in prayer.
The three prayers of the Mass ask God to arouse wills, turn hearts
and cure faults. These requests concern the people praying, including
people who have just received Communion. They are asking for changes
they cannot produce simply by deciding that their present condition
is satisfactory.

The difficulty is the apparent breadth of the promise. If its
confidence meant that every existing desire must be fulfilled,
the Secret's request for different desires would become hard to
understand. Aquinas explains prayer as raising the petitioner to
God, and Bede interprets the Markan instruction through faith
working by love and the duty to forgive. Together their arguments
explain how confident petition includes the petitioner's own
conversion.

\subsection{What it means to be made worthy}

The Epistle thanks the Father who has made the faithful worthy
of the saints' inheritance. Aquinas begins with the verb
\textit{fecit}, made. Worthiness is described as something
God produces in the recipients. It cannot therefore be used
as an independent qualification that explains why they first
received grace. The qualification itself belongs among the
gifts for which thanks are offered.\footnote{Aquinas,
\work{Super Colossenses} 1, lectio 3, on Col 1:12.}

Chrysostom gives the inference a concrete comparison. A king
can grant someone a governorship without making him capable
of exercising it. The honour may even expose the recipient
to ridicule if he is unfit. God's gift includes the fitness
as well as the inheritance. The recipients have not merely
been moved into a better location while remaining incapable
of the life given them.\footnote{Chrysostom,
\work{Homilies on Colossians} 2, on Col 1:12.}

The comparison explains why Paul's earlier demand to walk
worthily is compatible with gratitude for being made worthy.
The gift enables an ensuing life. Knowing God's will should
issue in good work; it is not a description of a recipient
who need never act. The Collect similarly asks for God's
action upon wills and for the faithful's pursuit of fruit
within one prayer.

Aquinas also explains the ``lot'' of the saints. An allotment
is received rather than seized by the beneficiary's own
choice. In this context it refers to God's judgment and
choice, not to a kingdom distributed by chance. Aquinas
identifies the saints' portion as God himself. This
identification clarifies why the Secret asks hearts to turn
toward God: the inheritance being sought concerns communion
with him.\footnote{Aquinas, \work{Super Colossenses} 1,
lectio 3, on \textit{in partem sortis sanctorum}.}

God's making fit does not make the person irrelevant to
what follows. The lesson names walking, bearing fruit,
growing and giving thanks. Aquinas observes that obedience
also disposes the mind to understand more deeply. The
petitioner asks for knowledge that becomes effective in
conduct and for conduct through which knowledge grows.
Prayer for a renewed will concerns that entire active life.

\subsection{Why pray for what God has promised?}

Jeremiah's promise raises a related question. If the Lord
already knows his purpose of peace and has promised return,
why does the passage speak of invocation and prayer?
Jerome explicitly asks why the people must pray if God is
to fulfil his own promise. He answers that invocation
and prayer prepare the captives to receive it. Their
petition participates in the promised restoration.\footnote{Jerome,
\work{In Ieremiam} V, on 29:10--14, PL 24, cols. 893--894.}

The chant's omissions can otherwise make its assurance
seem detached from the people addressed. The complete
letter requires whole-hearted seeking and rejects false
prophecy. The people must learn to want the return God
promises while living the life of fidelity he commands
where they are. The Introit's promise is thus consistent
with asking for aroused wills in the Collect. Both
concern recipients whom God's generosity summons to act.

Aquinas makes the same relationship explicit in his
explanation of Paul's praying and petitioning. Prayer
begins with the mind's ascent to God. A human petitioner
may first try to persuade another person to become
willing. In prayer to God, devotion and meditation
instead raise the petitioner toward him. The point
is to change the petitioner's relation to God, rather
than to overcome divine reluctance.\footnote{Aquinas,
\work{Super Colossenses} 1, lectio 3, opening explanation
of \textit{orantes et postulantes}.}

This account preserves a real request and a real gift.
Paul asks that the Colossians receive knowledge, virtue
and endurance; he does not merely describe an inward
exercise. Their prayer is one way in which they receive
and respond to God's good purpose. The gift can include
the changed capacity to desire and do what they request.

The Gradual's praise then concerns God's rescue rather
than the petitioner's self-sufficiency. Bellarmine reads
it as gratitude for past deliverance. The complete
Psalm has already rejected confidence in the speaker's
own bow and sword. Received help is a reason to thank
God, including while asking for help again.\footnote{Ps
43(44):4--10; Bellarmine, Psalm 43, explanation of the
corresponding native verses. Augustine's different account
of the perfect tense is developed in the preceding interpretation.}

\subsection{A cry that acknowledges both misery and mercy}

Psalm 129's depths describe the condition from which
petition is made. Bellarmine compares the sinner to
someone deep below the person to whom he cries.
God remains present by essence and power; the
distance concerns the sinner's dissimilarity to him.
The cry includes recognition of the need from which
the speaker seeks release.\footnote{Bellarmine,
\work{Commentary on the Book of Psalms}, Psalm 129,
explanation of vv. 1--2.}

Bellarmine then distinguishes two failures. Someone
who does not acknowledge his misery will not seek
medicine. Someone who does not know God's mercy
will despair of being cured. Penitence needs both
recognition and hope. The later verses of the
Psalm develop this reasoning through forgiveness
and trust in God's word; the two appointed
verses begin the actual appeal.\footnote{Bellarmine,
Psalm 129, explanation of vv. 4--5.}

Applied to the Collect and Postcommunion, that
distinction explains why asking for remedies need
not imply uncertainty about God's goodness.
The faithful seek medicine because they trust
the giver and acknowledge the defect. At the
Offertory, the opening cry is repeated. We can
understand its return as continued dependence
while gifts are brought, rather than as a claim
that offering has made dependence unnecessary.

The Secret explicitly includes the offerers
within the petition. God is asked to accept
offerings and prayers and to turn the hearts
of all. Fromage explains the prayer's passage
from earthly cravings toward heavenly desire.
The change concerns what the offerers seek,
not only whether their offering has been
received.\footnote{Fromage, vol. XI, pp. 512--513.}

This does not require abandonment of the
exiles' ordinary duties to build, plant and
care for their households. Those duties were
themselves commanded in the context of the
Introit. The distinction concerns possession
as something used responsibly and possession
as an attachment that resists God's call.
The Gospel's warning to leave a coat behind
provides a concrete case in which fidelity
requires the loss of an earthly good.

\subsection{Faith is tested by whom it trusts}

The Gospel also warns that a spectacular sign
can accompany a false claim. A person seeking
assurance may be attracted by the promise of
secret access to Christ. Jesus tells the
disciples beforehand how his coming differs.
Their trust must remain directed toward the
Lord who has spoken, rather than toward
whatever claim happens to promise immediate
relief.

Aquinas considers the disturbing qualification
that the false prophets would deceive even
the elect, if possible. He reports Origen's
explanation: a human being considered in
himself remains vulnerable to deception,
while divine election preserves its effect.
He also gives an alternative concerning
people elect only in appearance. These
are distinct accounts, not two descriptions
of the same person's final
state.\footnote{Aquinas, \work{Super Matthaeum}
24, lectio 3, on Mt 24:24, native anchor
87397. The Origen explanation is reported
by Aquinas, not cited here from an
independently examined Origen edition.}

Neither account makes apparent membership
a reason to disregard the warning.
The practical inference for this interpretation
is that a person praying confidently still
needs truthful discernment. A desire can
be intense without being a sound reason
to identify its promised fulfilment with
Christ. The Gospel's public lightning
and abiding words give the believer
reasons for refusing the hidden claimant.

The Trinity Preface identifies the God
to whom those desires are directed.
Its confession of the Father, Son and
Holy Spirit permits thanksgiving for
the Father's rescue and the Son's
redemption without treating their
action as competing favours.
Chrysostom expressly attributes the
whole saving benefit to both Father
and Son in explaining the
Epistle.\footnote{Chrysostom,
\work{Homilies on Colossians} 2,
on Col 1:13--14.}

Aquinas, commenting beyond the appointed
lesson, explains the Son as the image
of the invisible Father. He rejects
the idea that the Son is merely a
visible being of inferior nature.
The Son possesses the same invisible
divine nature, and all things are
created through him and sustained in
him. This account illuminates the
Preface's equal majesty. Heavenly
desire is directed to the living
God whom the prayer confesses,
with the angels who praise
him.\footnote{Aquinas,
\work{Super Colossenses} 1, lectio 4,
on Col 1:15--17. This is doctrinal
illumination of the Preface, not
part of the appointed Epistle.}

\subsection{Bede's explanation of the promise and forgiveness}

The Communion's assurance belongs to
Jesus' instruction about faith in
God. In Mark, the command to
forgive when standing in prayer
follows the promise about petition.
The whole instruction concerns
relations with God and with those
against whom the petitioner has
a grievance.\footnote{Mk 11:22--26.}

Bede begins his discussion of
mountain-moving with a real miracle
as a possibility when need requires.
He recounts Gregory of Neocaesarea's
prayer for a mountain to make room
for a church. The example matters
to his argument: he is not explaining
Jesus' promise by ruling out God's
extraordinary action in the
world.\footnote{Bede,
\work{In Marci evangelium expositio}
III, on Mk 11:22--23. The event
is related here as Bede's account.}

He also understands a spiritual
mountain as the devil in his pride.
In the next paragraph he identifies
perfect faith as faith working
through love. Its efficacy belongs
to a life transformed by God.
His account then addresses those
who have not reached such
perfection. They should seek
forgiveness, and he assures
them of an answer if they
first forgive those who have
sinned against them.\footnote{Bede,
\work{In Marci evangelium expositio}
III, on Mk 11:25. His adjacent
verse comments, rather than a
separate exposition of verse 24,
supply this explanation.}

The distinction prevents two
opposite mistakes. Extraordinary
power is not required before
an ordinary Christian may pray.
But the person asking pardon
cannot reserve the right to
withhold forgiveness as though
it were unrelated to the
petition. Bede grounds that
consequence in the next saying
of Jesus, not in an estimate
of the petitioner's emotional
intensity.

The promise therefore gives no
warrant for blaming someone
whose illness or danger
continues. Bede's argument
concerns faith, charity and
pardon; it supplies no
test by which one can
diagnose another person's
suffering as a failure to
believe. The Gospel of this
Mass has already described
grave danger among those
for whose sake the days
are shortened.

The Secret's request for
converted hearts can now
be related to the Communion
without weakening the assurance
of an answer. Forgiveness
may itself be part of
the change the petitioner
needs to receive. A
prayer concerning a
person one resents
includes submission of
that resentment to
God's judgment and
healing.

\subsection{The communicant's remaining faults}

The Postcommunion asks that
the received mysteries cure
what remains faulty in
the mind. Its confidence
concerns the medicinal
gift, while its petition
acknowledges a continuing
need. Fromage directly
explains this prayer
through the cure of
remaining defects.\footnote{Fromage,
vol. XI, p. 513.}

Augustine describes the
sacramental food and
drink as received in
memory of the Lord's
Passion for spiritual
health. God's invisible
operation brings about
its sanctifying
effect. That doctrinal
explanation concerns
more than a person's
determination to feel
better after reception.\footnote{Augustine,
\work{De Trinitate}
III.4.10.}

The resulting response
includes real effort:
knowledge should bear
fruit, the offerer
should accept conversion,
and the petitioner
should forgive.
Its possibility has
already been attributed
to God in the
Collect and Epistle.
The requested gift
includes a person
capable of sharing
the saints' inheritance
and praising the
Trinity with the
angels.

\begin{fourSenses}
\item[Literal.] Paul asks that
knowledge bear fruit and
thanks God for making
the faithful worthy.
Jesus' Markan instruction
places confident prayer
beside faith and
forgiveness. The three
orations ask God to
arouse, turn and heal
the faithful.
\item[Allegorical.] Redemption
through Christ's Blood
frees those brought
into his kingdom.
The Church offers
and receives in
Christ, confessing
Father, Son and
Holy Spirit in
one divine essence.
\item[Moral.] Ask for
good confidently,
seek God's will,
forgive those who
have sinned against
you, and submit
remaining faults
and disordered desires
to his cure.
\item[Anagogical.] The
saints' inheritance
and unending praise
fulfil the prayer
for heavenly desire.
The petitioner seeks
full healing and
communion with God.
\end{fourSenses}
